\documentclass[oneside,12pt]{article}          % ----- article

\input{../../../Latex/inputs/loading_packages.tex}			% ----- Packages utiles
\input{../../../Latex/inputs/preambule_exos.tex}         	% ----- Préambule feuilles de TD
\input{../../../Latex/inputs/mymacros.tex}                 % ----- Macros mathématiques
\input{../../../Latex/inputs/mise_en_page.tex} 		   	% ----- Macros de mise en page

\fancyhead[R]{2\up{nd} générale}            % En-tête droit
\fancyhead[C]{Résolution d'équations}
\fancyhf[CF]{\footnotesize{ \copyright ~ Tous droits réservés ~ \copyright}}                     % En-tête central
\setstretch{0.9}                           % Interligne

%%%%%%%%%%%%%%%%%%%%%%%%%%%%%%%%%%%%%%%%%%%%%%%%%%%%%%%%%%%%%%%%%%%%%%%%%%%%%%%%%%%%%%%%%%%%%%%%%%%%%%%%%%%%%
%%%%%%%%%%%%%%%%%%%%%%%%%%%%%%%%%%%%%%%%%%%%%%%%%%%%%%%%%%%%%%%%%%%%%%%%%%%%%%%%%%%%%%%%%%%%%%%%%%%%%%%%%%%%%
%%%%%%%%%%%%%%%%%%%%%%%%%%%%%%%%%%%%%%%%%%%%%%%%%%%%%%%%%%%%%%%%%%%%%%%%%%%%%%%%%%%%%%%%%%%%%%%%%%%%%%%%%%%%%


\begin{document}

\section{Équations du premier degré}

\subsection{Échauffement}
Résoudre dans $\RR$ chacune des équations suivantes :

\renewcommand{\arraystretch}{1.5}
\begin{tabular}{p{0.3\linewidth} p{0.3\linewidth} p{0.3\linewidth}}
	a) $3x+15=0$ & b) $2x-11=7x+4$ & c) $7x-20=3x+9$
\end{tabular}



\subsection{Résolutions différenciées}
Les équations suivantes admettent-elles des solutions dans $\ZZ$ ? Dans $\DD$ ? Et dans $\QQ$ ?
	
\renewcommand{\arraystretch}{1.5}
\begin{tabular}{p{0.3\linewidth} p{0.3\linewidth} p{0.3\linewidth}}
	a) $12x + 5 = 9x + 12$ & b) $-4x + 1 = 6x - 18$ & c) $19x - 8 = 13x + 14$
\end{tabular}



\subsection{Équations avec fractions}
Résoudre dans $\RR$ les équations suivantes :
	
\renewcommand{\arraystretch}{1.5}
\begin{tabular}{p{0.3\linewidth} p{0.3\linewidth} p{0.3\linewidth}}
	a) $\dfrac{x}{3} + 1 = 2x + 2$ &
	b) $\dfrac{x}{2} + \dfrac{1}{3} = \dfrac{x}{3} - \dfrac{1}{2}$ &
	c) $\dfrac{3x}{2} - \dfrac{5}{6} = 2 + \dfrac{2x}{3}$ 
\end{tabular}



\subsection{Plus difficile}
Résoudre dans $\RR$ les équations suivantes :
	
\renewcommand{\arraystretch}{1.5}
\begin{tabular}{p{0.3\linewidth} p{0.3\linewidth} p{0.3\linewidth}}
	a) $\dfrac{2x - 1}{5} = x + \dfrac{3}{4}$ &
	b) $\dfrac{3x}{4} + \dfrac{2}{3} = \dfrac{x - 1}{2}$ &
	c) $\dfrac{6x}{11} - \dfrac{x}{3} + 1 = \dfrac{2x + 21}{4}$
\end{tabular}



\subsection{Du faux second degré}
Résoudre dans $\RR$ chacune des équations suivantes :

\begin{tabular}{p{0.3\linewidth} p{0.3\linewidth} p{0.3\linewidth}}
	a) $8x + x^2 - 15 = x^2 + 4x - 1$ &
	b) $x(x + 3) = x^2 - 1$ &
	c) $2x^2 + 5x - 9 = 2x(x + 8)$
\end{tabular}



\section{Autour des valeurs interdites}

\subsection{Rappel nécessaire}
On considère les trois fonctions $f$ , $g$ et $h$ définies ci-après :

\begin{tabular}{p{0.3\linewidth} p{0.3\linewidth} p{0.3\linewidth}}
	$f : x \mapsto \dfrac{x + 1}{2x - 4}$ &
	$g : x \mapsto \dfrac{3x + x^3}{-2x + 9}$ &
	$h : x \mapsto \dfrac{x^2 + x - 3 }{3x + 10}$
\end{tabular}

\begin{enumerate}
	\item Déterminer la valeur interdite de chacune de ces fonctions.
	\item Donner alors leur domaine de définition.
\end{enumerate}



\subsection{Dénominateur variable}
Pour chacune des équations ci-après : déterminer leur valeur interdite, puis les résoudre dans $\RR$ :

\begin{tabular}{p{0.3\linewidth} p{0.3\linewidth} p{0.3\linewidth}}
	a) $\dfrac{2}{x - 3} = 1$ &
	b) $\dfrac{x + 3}{x} = 2$ &
	c) $\dfrac{1}{x - 2} = 3$
\end{tabular}



\subsection{Avec deux dénominateurs}
Pour chacune des équations ci-après : déterminer leurs deux valeurs interdites, puis les résoudre dans $\RR$

\begin{tabular}{p{0.3\linewidth} p{0.3\linewidth} p{0.3\linewidth}}
	a) $\dfrac{1}{x} + \dfrac{1}{2x + 1} = 0$ &
	b) $\dfrac{2}{x+1} - \dfrac{1}{x} = 0$ &
	c) $\dfrac{3}{x} + \dfrac{1}{x + 1} = 1$
\end{tabular}





\section{Équations produit nul}

\subsection{Première application}
Résoudre dans $\RR$ chacune des équations suivantes :

\begin{tabular}{p{0.32\linewidth} p{0.32\linewidth} p{0.32\linewidth}}
	a) $(3x+2)(4x-5)=0$ & b) $(6x-2)(2x+6)=0$ & c) $\left(\frac{3}{2}x-2\right)(8-6x)=0$ \\
	d) $(2x+1)(4x-3)(6x+18)=0$ & e) $\left(\frac{x}{3}+2\right)\left(2x-\frac{1}{3}\right)(7-4x)=0$ & f) $\left(\frac{2x}{3}+4\right)(3x-12)\left(3+\frac{x}{2}\right)=0$
\end{tabular}



\subsection{Avec une factorisation}
Résoudre dans $\RR$ chacune des équations suivantes à l'aide d'une factorisation :

\begin{tabular}{p{0.45\linewidth} p{0.45\linewidth}}
	a) $x(3x+2)-6(3x+2)=0$ & b) $4x^2-25=0$ \\
	c) $9x^2+12x+4=0$ & d) $(3x-4)^2-(x+5)^2=0$ \\
	e) $3x^2-27x=0$ & f) $(2x+6)(4-x)-(3-2x)(2x+6)=0$ \\
	g) $3x^2-17x+25=3x-x^2$ & h) $(7x-8)(4x+20)=(3x-1)(4x+20)$ \\
	i) $(2x-5)^2=x^2$ & j) $5x^2-17x=6x+x^2$
\end{tabular}



\subsection{Avec des factorisations}
Résoudre dans $\RR$ chacune des équations suivantes à l'aide d'une ou plusieurs factorisations :

\begin{tabular}{p{0.45\linewidth} p{0.45\linewidth}}
	a) $49(2x+1)-x^2(2x+1)=0$ & b) $(2x^2+1)^2-(2x^2+4x)^2=0$ \\
	c) $16x^4-81=0$ & d) $(4x-5)^2-5x+4x^2=0$ \\
	e) $(x^2+3x-1)^2-(x^2-3x+1)^2=0$ & f)$(x^2-16)(x^2+3x)-(3x+9)(x^2-16)=0$
\end{tabular}



\subsection{En deux étapes}
\begin{enumerate}
	\item \begin{enumerate}
		\item Développer l'expression $(2x-1)(x+3)$
		\item Résoudre dans $\RR$ l'équation $2x^2+5x-3=0$
	\end{enumerate}
	\item \begin{enumerate}
		\item Développer l'expression $x(3x-2)(x+2)$
		\item Résoudre dans $\RR$ l'équation $3x^3+4x^2-4x=0$
	\end{enumerate}
	\item \begin{enumerate}
		\item Développer l'expression $(x-2)(2x+4)(5-3x)$
		\item Résoudre dans $\RR$ l'équation $3x^3-5x^2-12x+20=0$
	\end{enumerate}
\end{enumerate}



\section{Résolution de problèmes}




\end{document}